\documentclass[11pt]{amsart}

\usepackage[T1]{fontenc}
\usepackage[utf8]{inputenc}
\usepackage{amsmath,amssymb,amsthm}
\usepackage[margin=1in]{geometry}
\usepackage{hyperref}

\newcommand{\LL}{\Lambda}
\newcommand{\OO}{\mathbb{O}}
\newcommand{\ZZ}{\mathbb{Z}}
\newcommand{\FF}{\mathbb{F}}
\newcommand{\RR}{\mathbb{R}}
\newcommand{\Nrm}{N}

\title[Addendum (2026-05-24): the $L/2L$ parallel]%
{Addendum to\\
\emph{Verification and Commentary: The Baez--Egan Embedding of the
Leech Lattice in the Exceptional Jordan Algebra}\\
\large The $L/2L$ parallel and how our work builds on it}

\author{Claude Opus 4.7}
\address{Anthropic, at the direction of Jens K\"oplinger}
\date{24 May 2026}

\begin{document}
\maketitle

\begin{abstract}
This is a dated addendum to the original 16 April 2026 Verification
and Commentary on the Baez--Egan embedding
(\texttt{paper/Baez-Egan\_Leech\_Verification.pdf}, 8\,pp).  Re-reading
that work in light of our subsequent paper \emph{An order on the Leech
lattice from a $\ZZ_3$-symmetric triple-octonion product} (v4,
24 May 2026), we notice an explicit lattice-algebraic parallel
between Baez/Egan's ``key lemma'' (Part 9, $L_1 \cap L_2 = 2 L_0$)
and Wilson's identity ($L\bar s \cap Ls = 2 L$), and we identify
exactly how Section~A.1 of our paper (\emph{the mod-2-quotient
reading}) builds on top of that shared identity.  The purpose of
this addendum is to credit the shared lattice-algebraic foundation
correctly and to delineate what is new in our use of it.
\end{abstract}

% ===================================================================
\section{The shared key lemma}
% ===================================================================

Both Wilson and Baez/Egan rest on the same lattice-algebraic
identity, in different notations:

\begin{center}
\begin{tabular}{l@{\hskip 2em}l}
Baez/Egan (Part 9, Section~4 of original PDF): &
$L_1 \cap L_2 = 2 L_0$, \\[2pt]
Wilson (2009, Section~2): &
$L\bar s \cap Ls = 2 L$. \\
\end{tabular}
\end{center}

In Baez/Egan's setup, $L_0 \cong E_8$ is self-dual; $L_1, L_2 \cong
\sqrt{2}\,E_8$ are sublattices of~$L_0$ arising from a common
rotation~$R$, and the duality argument $\mathrm{dual}(L_1 \cap L_2) =
\tfrac12 L_1 + \tfrac12 L_2 = \tfrac12 L_0$ gives
$L_1 \cap L_2 = 2 L_0$ (Egan).

In Wilson's setup, $L = D_8^+$ is the octonion-realised $E_8$; the
sublattices~$L\bar s$ and~$Ls$ are the images of left-multiplication
by $\bar s$ and~$s$ respectively (so each is isomorphic to
$\sqrt{2}\,E_8$, since $\Nrm(s) = \Nrm(\bar s) = 2$), and Wilson
states the identity directly.

The two identities are the same algebraic fact about $E_8$ and a
pair of $\sqrt{2}\,E_8$ sublattices: Wilson's version is the
octonion-specific instance; Baez/Egan's is the geometric stripping
of it (Egan's Part 9).  Either form was already in the literature
before this project began.

% ===================================================================
\section{What Baez/Egan use it for, and what they did not do}
% ===================================================================

Baez/Egan use the key lemma in Part~9 as a step in establishing
that the lattice
\[
  L_L \;=\; \{(a, b, c) \in L_0^3 :
    a, b, c \in L_0;\;
    a{+}b, a{+}c, b{+}c \in L_1;\;
    a{+}b{+}c \in L_2\}
\]
is isometric to $\sqrt{2}\,\LL$.  The identity $L_1 \cap L_2 = 2 L_0$
is used in showing both (i)~$L_L$ is even and self-dual, and (ii)~$L_L$
has no vector of squared norm~$<4$.  This is the use in Egan's
Part~9 of the original $n$-Category Caf\'e series.

Baez/Egan do \emph{not}, anywhere in the Part~9 / Part~10 sequence,
descend the octonion product on $L_0 = E_8$ modulo~$2 L_0$ to read
$L_0 / 2 L_0$ as an $\FF_2$-bilinear algebra and analyse the sublattice
images $L_1 / 2 L_0$, $L_2 / 2 L_0$ in that algebra.  The key lemma
is used at the level of lattices in $\OO^3$, not at the level of an
$\FF_2$-algebra structure inside $L_0 / 2 L_0$.

% ===================================================================
\section{What our Section A.1 builds on the same identity}
% ===================================================================

Section~A.1 of our paper (\emph{The mod-2-quotient reading of
Tables~A.2 and~A.3}, v4) starts from the same identity, in Wilson's
form:
\begin{equation}\label{eq:key}
  L\bar s \cap Ls \;=\; 2 L, \tag{$\dagger$}
\end{equation}
and uses it as the foundation for an explicitly algebraic reading.
The steps:

\begin{enumerate}
\item From~\eqref{eq:key}, $2 L \subseteq L\bar s$ and $2 L \subseteq Ls$;
  applying~$\sigma$ (which fixes $2 L$, since it is a coordinate
  permutation of an even integer lattice),
  $2 L \subseteq \sigma(L\bar s)$ and $2 L \subseteq \sigma(Ls)$.
\item $L \cdot L \subseteq L$ (Coxeter; classical) and
  $\ZZ$-bilinearity descend the octonion product to an
  $\FF_2$-bilinear product $\bar\mu$ on
  $\bar L := L / 2 L \cong \FF_2^{\,8}$.
\item The sublattices $\sigma(L\bar s)$ and $\sigma(Ls)$ are full
  preimages of $4$-dimensional $\FF_2$-subspaces
  \(V := \sigma(L\bar s) / 2 L\) and \(W := \sigma(Ls) / 2 L\),
  and from $\sigma(L\bar s) + \sigma(Ls) = L$ together
  with~\eqref{eq:key} (applied after~$\sigma$),
  $\bar L = V \oplus W$.
\item The closure statements of Lemmas~4.3 and~4.4 in our paper
  become, mod~$2L$:
  \(V\) is a (two-sided) ideal of $\bar L$;
  \(W\) is a subalgebra of $\bar L$.
\item Under the natural plus-type quadratic form $q(v + 2L) :=
  \Nrm(v) / 2 \pmod 2$ on $\bar L$, direct verification shows that
  both~$V$ and~$W$ are totally isotropic, hence Lagrangians of
  $(\bar L, q)$; the decomposition $\bar L = V \oplus W$ is a
  Witt decomposition into a complementary pair of Lagrangians.
\end{enumerate}

Step~(1) uses identity~\eqref{eq:key} directly.  Steps~(2)--(5) are
new in our paper: they descend the octonion product to an
$\FF_2$-bilinear product on $L / 2 L$, read the sublattices as
ideal/subalgebra in that descended algebra, and recognise the whole
picture as a polarization of an $\FF_2$-orthogonal space.  Baez/Egan
in Parts~9--10 do not take this $\FF_2$-algebra step.

% ===================================================================
\section{Attribution and overlap}
% ===================================================================

What is shared with Baez/Egan:

\begin{itemize}
\item The lattice-algebraic identity $L_1 \cap L_2 = 2 L_0$ (Egan's
  Part~9 key lemma), equivalent to Wilson's $L\bar s \cap Ls = 2 L$.
  Egan's duality-based proof of this is in the original PDF~\S4 of
  the 16~April 2026 Verification and Commentary, and is reproduced
  here for the reader's convenience.
\item The role of this identity in setting up a Leech-lattice
  construction inside $L_0^3$ via three sum-of-components conditions
  (Wilson 2009 in octonion language; Egan Part~9 in lattice
  language).
\end{itemize}

What is not in Baez/Egan and is new in our paper:

\begin{itemize}
\item The descent of the octonion product modulo $2 L$ to an
  $\FF_2$-bilinear product on $L / 2L$.
\item The identification of $V$ and $W$ as ideal and subalgebra of
  this descended algebra.
\item The Witt-decomposition (polarization) reading of $L / 2L$ into
  complementary Lagrangians under the natural plus-type quadratic
  form.
\end{itemize}

The main paper (v4, Section~A.1) cites Wilson~2009 for the
sublattice identities and notes Baez/Egan in~\S6 as part of the
related-work discussion.  The cross-credit recorded here -- that
our Section~A.1 builds on the same lattice-algebraic key lemma
that Egan's Part~9 stripped to its geometric core -- is implicit
in the main paper's references but is made explicit in this
addendum.

% ===================================================================
\section{What it would take to close the gap further}
% ===================================================================

A natural follow-up question: does Baez/Egan's framing extend to the
$\FF_2$-algebra picture of Section~A.1?  In their setup the
generalised lattice $L_L \subset L_0^3$ is constructed without
explicit octonion multiplications; once those multiplications are
brought back in (as in Wilson 2009, and as required for our
$\bar\mu$ descent), the $\FF_2$-algebra picture is available.
Egan's Part~9 chooses to abstract away from octonion multiplication;
that abstraction is what makes our $\FF_2$-algebra reading absent
from his account.  In this sense, the present work is a continuation
of the Wilson-side of the Wilson--Egan parallel, not the Egan-side.

Conversely, the Part~10 construction (lifting $\LL$ to the
$27$-dimensional Jordan-algebra ambient $\widetilde{\LL} \subset
\mathfrak{h}_3(\OO)$, closed under $2(XY + YX)$) is a different
direction: it embeds $\LL$ into a $27$-dimensional ring,
whereas our mod-2 picture works within $24$-dimensional $L / 2L$
inherited from~$L$ alone.  The two routes are not in competition;
they are different uses of the same shared key lemma.

% ===================================================================
\section*{References}
% ===================================================================

\begin{itemize}
\item Original Verification PDF:
  \texttt{paper/Baez-Egan\_Leech\_Verification.pdf}, 16 April 2026
  (Claude, Anthropic).
\item Our main paper:
  \emph{An order on the Leech lattice from a $\ZZ_3$-symmetric
  triple-octonion product}, v4 (24 May 2026), Section~A.1
  ``The mod-2-quotient reading of Tables A.2 and A.3''.
\item Wilson, R.~A., \emph{Octonions and the Leech lattice},
  J.~Algebra \textbf{322} (2009), 2186--2190.
\item Baez, J.~C.\ (with G.~Egan and others),
  \emph{Integral octonions} (parts 9 and 10),
  $n$-Category Caf\'e, November--December 2014.
\end{itemize}

\end{document}
